\chapter{Background and Related Work}\label{ch:background}

This chapter provides the technical and regulatory context for the project. It
opens with the two standards and the security problems they respond to each other. Compliance
mapping as an activity takes place, then the NLP and LLM techniques the project
builds on, and the chapter closes on prior studies of automated compliance
results and analysis.

\section{Background}

\subsection{IoT security and \standardA{}}

Consumer IoT devices are network connected by design, often built on
limited hardware, and frequently shipped with weak defaults. Recurring
weaknesses include poor authentication mechanisms, unprotected software
updates, unencrypted communication, unnecessarily exposure of interfaces, and
inadequate protection of personal data~\cite{sicari2015security,vardakis2024smarthome}. Because such devices are used in homes and workplaces, even minor security failures can lead into concrete
privacy and safety risks~\cite{mocrii2018smarthomes}.

\standardA{} (\textit{Cyber Security for Consumer Internet of Things: Baseline
Requirements}) addresses these weaknesses with thirteen groups of provisions,
including no universal default passwords, vulnerability disclosure, keeping
software updated, secure communication, minimising attack surfaces, and
ensuring the protection of personal data~\cite{etsi303645}. The standard is
outcome-focused. It states what a device must achieve rather than prescribing
implementations, and it segregate mandatory provisions (``shall'') from
recommendations (``should''). Its provisions are managed in numbered clauses
(e.g., provision~5.3-3 on update mechanisms), which makes it a practical anchor
for systematic mapping.

\subsection{AI system security and \standardB{}}

AI systems face threats that classical software security does not fully cover:
training data poisoning, model manipulation and theft, adversarial inputs, and
prompt injection in systems built on large models~\cite{chen2024mlsecurity,
he2025aisecurity}. Security work for AI must therefore address the model
life cycle (data collection, training, deployment, maintenance, and
disposal) and not only the running system.

\standardB{} (\textit{Securing Artificial Intelligence: Baseline Cyber
Security Requirements}) structures its provisions along exactly that life
cycle: secure design (clause 5.1), secure development (5.2), secure deployment
(5.3), secure maintenance (5.4), and secure end of life (5.5)~\cite{etsi304223}.
Within these phases it covers AI-specific control areas such as threat and risk
management, documentation and auditability of models, protection of datasets
and prompts, supply-chain security, security testing, behavioural monitoring,
and secure disposal of data and models. The provision-based, life-cycle
structure parallels the clause organisation of \standardA{}, which is what
makes a pairwise mapping between the two documents feasible in the first
position.

\subsection{Compliance mapping}

Compliance mapping is the activity of relating the requirements of one
framework to those of another, so that evidence collected for one can be
used again for the other and gaps become visible. Studies of security standards
report that overlapping controls, inconsistent terminology, and fragmented
requirements are common across frameworks, and that mapping them requires
interpretation of purpose and security objective, not just wording
~\cite{djebbar2023comparative,mussmann2020mappings}. Similar outcomes come
from the GDPR domain, where manual level compliance checking is described as
time consuming and dependent on scarce legal technical
expertise~\cite{cejas2023gdpr,negri2024gdpr}. Castellanos Ardila et
al.~\cite{castellanos2022compliance} finds that manual compliance
checking of software processes is labour intensive and hard to scale.

Automated approaches try to minimize this effort by withdrawing requirements from
the documents and comparing them algorithmically, classifying pairs as
identical, related, partially overlapping, or unrelated~\cite{sleimi2021legal}.
Beyond saving effort, automation promises consistency: the same criteria are
applied to every pair, which manual mapping cannot guarantee. The open question
is accuracy, and that is the subject of this thesis.

\subsection{NLP and LLM techniques for requirement analysis}

\paragraph{Semantic similarity.}
Standards frequently express similar objectives with different terminology"
device-focused" language in \standardA{}, "life-cycle and governance" language in
\standardB{}. Keyword matching fails on such vocabulary mismatch, a problem
long recognised in requirements traceability
research~\cite{mahmoud2015semantics,rajpathak2022linking}. Semantic approaches
map text into vector spaces where related meanings land close together, and
have been shown to outperform lexical matching for requirement tracing
~\cite{mahmoud2015semantics}.

\paragraph{NLP pipelines.}
NLP has been applied across requirements engineering tasks where extracting domain
concepts, detecting language defects, and establishing traceability links
~\cite{zhao2021nlp4re,sonbol2022representation} takes place. A typical pipeline extracts
normative statements, normalises the text, and converts each requirement into
a numerical representation suitable for comparison. Legal and regulatory text
adds difficulties of its own: context-dependent terminology, implicit
meaning, and complex cross-references~\cite{ariai2025legal}.

\paragraph{Transformer models.}
Transformer architectures learn contextual representations of language through
self-attention. BERT~\cite{devlin2019bert} provides bidirectional contextual
embeddings that transfer well to many NLP tasks. For sentence-level
comparison, Sentence-BERT~\cite{reimers2019sbert} fine-tunes a Siamese
architecture so that cosine similarity between sentence embeddings reflects
semantic relatedness, making large-scale pairwise comparison cheap.
SecureBERT~\cite{aghaei2022securebert} adapts a RoBERTa-based model to
cybersecurity text, motivating its inclusion here as a domain-tuned
alternative. One caveat known from the literature is that the embedding spaces
of contextual models are anisotropic: vectors concentrate in a narrow cone,
so unrelated sentences can still receive high cosine
similarity~\cite{ethayarajh2019contextual}. This property matters for the
results in Chapter~\ref{ch:results}.

\paragraph{Large language models.}
LLMs can be prompted to compare two requirements directly and output a
relationship judgement, capturing nuances that similarity scores flatten.
Reported drawbacks include computational cost, limited reproducibility, and
sensitivity to prompt design~\cite{jamshidi2025llmiot,rodriguez2024llm}.

\paragraph{Gap analysis.}
Mapping two standards also serves gap analysis: identifying objectives covered
by one document but not the other. \standardA{} concentrates on device-level
baseline security while \standardB{} concentrates on the AI life cycle, so
their relationship is expected to consist mostly of partial overlaps and
complementary provisions rather than one-to-one equivalences
~\cite{etsi303645,etsi304223}. The task definition therefore covers partial and
absent alignment as well as strong matches.

\section{Related work}

\subsection{Manual comparative mapping of security standards}\label{sec:manualmapping}

The closest published work to this thesis is manual. Djebbar and
Nordstr\"{o}m~\cite{djebbar2023comparative} compare ISO/IEC 27001:2022,
ISA/IEC 62443-3-3:2019 and \standardA{}, and their method is explicitly
interpretive: controls are matched by \emph{teleological interpretation}, that
is, by the goal a control is intended to realise rather than by the words it
uses, and the authors state a deliberate preference for under-mapping over
over-mapping when a match is uncertain. Their taxonomy is binary: two
controls are \emph{Equivalent} when wording and purpose coincide, and
\emph{Related} when the wording differs but the security objective is the same.

Their results are the reference point for what a careful manual mapping of
\standardA{} produces. All 68 provisions of \standardA{} could be aligned with
29 controls of ISO/IEC 27001:2022, while 64 of that standard's 93 controls had
no counterpart in \standardA{} at all, almost entirely the organizational
and people controls, which have no meaning for a device. For ISA/IEC 62443-3-3
they report 84 of 100 requirements mapped, and tabulate the unmapped remainder
individually rather than reporting only a coverage figure. The complete
mappings are given as appendix tables, one row per aligned control pair.

Two features of that study shape this one. The first is the sparsity: a
standards pair drawn from different domains produces far more non-matches than
matches, and a mapping method that does not reproduce that asymmetry is not
reproducing the task. The second is the form of the output: a typed
relation, a coverage statement, and an explicit list of what was left unmapped.
Section~\ref{sec:coverage} produces those three artefacts automatically.

\subsection{Automated compliance mapping in cybersecurity}

Bar-Haim et al.~\cite{barhaim2023cloud} automate the assessment of
organizational cybersecurity posture in cloud environments by mapping
free-text security requirements to two control frameworks: a NIST-based
framework of 280 controls in 22 families, and a finer-grained three-level risk
framework of 169 categories. They formulate mapping as supervised text-pair
classification with RoBERTa and SBERT models and use hard negative sampling to
force the models to separate semantically close but non-equivalent controls.
Their best combined model reaches about 72\% precision@1 for control-family
mapping and 55\% for individual controls, degrading as granularity increases.
The authors conclude that the approach suits a human-in-the-loop workflow
rather than autonomous operation, a conclusion this thesis revisits for the
standard-to-standard setting.

\subsection{NLP and semantic techniques for requirement mapping}

Sapkota et al.~\cite{sapkota2016regcmantic} present RegCMantic, a framework
that maps regulatory guidelines to organizational processes. Regulatory
documents are parsed into a uniform XML structure, requirement entities
(subject, obligation, action) are extracted into an ontology, and mapping uses
three similarity measures over topics, core entities, and auxiliary entities.
In a pharmaceutical case study the extended framework improves extraction
F-measure from 0.83 to 0.91 and performs mapping evaluated against
expert-created references at a 0.85 similarity threshold. The study
demonstrates that ontology-supported NLP can automate requirement mapping,
while noting that ambiguity and implicit, domain-specific meaning remain the
hard part. The same difficulties appear with the standards studied here.

\subsection{LLM-based compliance analysis}

Rodríguez et al.~\cite{rodriguez2024llm} investigate LLMs for large-scale
privacy policy analysis, comparing ChatGPT and Llama~2 against earlier
symbolic and machine-learning baselines on the MAPP and OPP-115 datasets.
With a two-shot prompt design their best configuration exceeds 93\% F1,
outperforming prior techniques while requiring less annotated data and
engineering effort. The study is limited to privacy policies, but it provides
strong evidence that prompted LLMs can support compliance-oriented text
analysis, motivating the inclusion of an LLM classification method in this
project.

\subsection{Positioning of this thesis}

Prior work automates mapping between requirements and control frameworks
~\cite{barhaim2023cloud}, between regulations and business processes
~\cite{sapkota2016regcmantic}, and classification within privacy policies
~\cite{rodriguez2024llm}. Mapping two security standards to each other, with
a typed relationship taxonomy instead of a binary match, has not to our
knowledge been evaluated systematically for the IoT/AI standard pair.

The relationship to \cite{djebbar2023comparative} is the one worth stating
plainly, because this thesis attempts by machine what that study did by hand.
It takes the same interpretive task to a standard pair that study does not
cover, \standardA{} against \standardB{}; it refines the Equivalent/Related
distinction into six typed relations, so that partial and directional
alignment can be expressed instead of collapsed; and it asks a question a manual
study cannot answer about itself, namely how much of that judgement survives
automation and at what error rate. Where the manual study reports its mapping,
this one reports a mapping together with a measurement of how often it is wrong,
scored against an annotated reference set on which the trivial answer is held to
a known score. The comparison of eleven methods that follows measures how closely
each approximates the judgement \cite{djebbar2023comparative} exercised
directly.
